\documentclass[../main.tex]{subfiles}

\begin{document}

\chapter{Foundations and Intuitions}

This note covers basic concepts as well as intuitions through
the introduction of Large Language Models.

\section{Fundamental Terms and Concepts}

This section is an introduction to basic terms and concepts that
will be used throughout my notes. That being said, more advanced
concepts will be introduced in latter notes. Please feel free
to skip over this section if you already know the bold words.

\begin{definition}{}{}
\textbf{Machine learning} is a subset of artificial intelligence
that learns from data. \textbf{Deep learning} is a subset of
machine learning that utilizes neural networks.
\end{definition}

The primary purpose of ML is to make AI ``smarter''. For instance,
calculators are ``dumb'' in a sense that it requires an input and
can only do certain tasks such as outputting calculated input.
However, we can make it ``smarter'' by teaching data. We can clearly
see that WolframAlpha is smarter than a normal scientific calculator.
(Please note that this is just an analogy not an example because
WolframAlpha is a knowledge engine, not a machine learning model.
But I mean you get the point.)

On the high level, you can think of an AI model as a giant function
with billions of array of numbers, called parameters.

\begin{definition}{}{}
\textbf{Parameters} are the internal values in a model composed
with \textbf{weights} and \textbf{biases} that are learned by the
model. After the parameters are finalized with the training phase,
the models generate \textbf{predictions} through the process called
\textbf{inference}.
\end{definition}

Indeed, models are prediction machines based on the trained
parameters. One fun fact here is that it generally costs way less
to run inference than to train the model. More in-depth explanation
will be available in latter notes where we discuss neural nets.

\subsection{ML Paradigms}

In ML, there are three primary learning paradigms: supervised
learning, unsupervised learning, and reinforcement learning.
First, what does it mean to ``learn'' and ``train''?

\begin{definition}{}{}
Machine learning utilizes the process of \textbf{training} to
learn from data with algorithms. A common example would be
dog and cat recognition from picture with \textbf{labeled} data.
\end{definition}

There are multiple ways in which models are trained. One of the
most used method out of the three paradigms is supervised learning.

\begin{definition}{}{}
When a model is trained with \textbf{supervised learning},
the model utilizes labeled dataset, or examples of input-output
pairs, to learn.
\end{definition}

With supervised learning, there are two primary predictions that
models make.

\begin{definition}{}{}
\textbf{Classification} is a prediction where the label of an
input data is being predicted. \textbf{Regression} is utilized
to predict continuous values.
\end{definition}

Classic examples of classification would be categorizing dog/cat
pictures and email spam filtering. An example of regression would
be the house price over years.

The second most common learning paradigm is unsupervised learning.

\begin{definition}{}{}
\textbf{Unsupervised learning} is a process of training a model
with unlabeled data.
\end{definition}

How is it possible for a model to learn without labeled data?
One primary example algorithm is $k$-means clustering,
where the model categorizes the data with relevance and similarity.
To cluster high dimensional data, we often perform
\textbf{dimensionality reduction} to reduce the dimension for
clustering with the trust in \textbf{manifold hypothesis}.
I won't go to deep into unsupervised learning and $k$-means clustering
in this note, but you can check out this
\href{https://www.naftaliharris.com/blog/visualizing-k-means-clustering/}
{awesome visualization} of $k$-means clustering.

The third most used learning paradigm is reinforcement learning.

\begin{definition}{}{}
\textbf{Reinforcement learning} is a process of training models
through interactions with environment and reward system.
\end{definition}

Personally, reinforcement learning is more complicated compared to
the two paradigms introduced above. On the high level, an agent will
take an action in an environment, where it returns \textbf{rewards}
based on the action. The agent will then update its policies with
the provided reward. The process is continued until the ``adequate''
amount is reached.

Training a model can be understood as writing a best fit
mathematical function to training datasets that captures the
holistic trend. There are few terms that describe the function.

\begin{definition}{}{}
A model is \textbf{overfit} if it learns too much from the data
that it does not effectively capture that holistic pattern but
small \textbf{noises}. On the other hand, an \textbf{underfitting}
model is formed when it does not learn effectively from the data.
\end{definition}

On the high level, our goal is to generalize and capture the general
pattern and prevent overfitting or underfitting. We normally would
stop training a model when it does not appear to improve, or is
degrading, but how do we assess it?

\begin{definition}{}{}
To prevent overfitting and to evaluate the model during training,
\textbf{validation} is utilized with \textbf{unseen data}.
Unlike validation, \textbf{test sets} are used at the end of the
training to evaluate the finalized model.
\end{definition}

We have discussed the fundamental terms and concepts that would
repeatedly appear in the following notes. Now let's build our
high level intuitions on ML with the introduction to LLMs!

\end{document}


